% ============================================================
%  Juros e Desconto — Apostila Premium | PratiqueJá
%  Engine: XeLaTeX
% ============================================================
\documentclass[12pt]{article}
\usepackage{../../pratiqueja}

\newcommand{\hlm}[1]{{\color{darkblue}#1}}

\setsubject{Juros e Desconto}

\begin{document}
\pagenumbering{arabic}
\setstretch{1.2}
\color{bodytext}

\heroheader{Juros e Desconto}%
           {Regimes simples e compostos, taxas e descontos}%
           {Matemática · Avançado}

\setlength{\columnsep}{18pt}
\setlength{\columnseprule}{0.5pt}
\def\columnseprulecolor{\color{exborder}}

\begin{multicols}{2}

\TW{JS}{darkblue}{Juros Simples}{Juros calculados apenas sobre o capital inicial}

\regra{
Nos \hl{juros simples}, o valor dos juros é calculado \textbf{apenas sobre o capital inicial} durante todo o período. A base de cálculo não muda — os juros não geram novos juros.
}

\formula{$J = C \cdot i \cdot t$}

\begin{itemize}
  \item $J$ — juros simples
  \item $C$ — capital inicial (valor principal)
  \item $i$ — taxa de juros por período (em decimal)
  \item $t$ — tempo (mesmo período da taxa)
\end{itemize}

\SH{Montante}

O montante $M$ é capital mais juros:

\formula{$M = C + J = C \cdot (1 + i \cdot t)$}

\exemp{
\textbf{R\$ 1.000,00 a 5\% a.m. por 6 meses:}

$J = 1000 \cdot 0{,}05 \cdot 6 = 300$

$M = 1000 + 300$

$\hlm{M = \text{R\$ } 1.300{,}00}$
}

\TW{JC}{darkblue}{Juros Compostos}{Juros sobre capital + juros acumulados}

\regra{
Nos \hl{juros compostos}, a cada período o montante (capital + juros acumulados) torna-se o novo capital. Os \textbf{juros geram mais juros} — o crescimento é exponencial.
}

\formula{$M = C \cdot (1 + i)^t$}

\begin{itemize}
  \item $M$ — montante final
  \item $C$ — capital inicial
  \item $i$ — taxa de juros por período (em decimal)
  \item $t$ — número de períodos
\end{itemize}

O juros composto total: $J = M - C$

\exemp{
\textbf{R\$ 1.000,00 a 10\% a.a. por 2 anos:}

$M = 1000 \cdot (1{,}10)^2 = 1000 \cdot 1{,}21$

$\hlm{M = \text{R\$ } 1.210{,}00}$
}

\TW{JS×JC}{badgepurple}{Comparativo: Simples × Composto}{Quando usar cada regime}

\regra{
Para $t = 1$ período, ambos os regimes produzem o mesmo resultado. Para $t > 1$, os juros compostos sempre superam os simples.

\begin{itemize}
  \item \textbf{Juros Simples} — crescimento linear; operações de curto prazo
  \item \textbf{Juros Compostos} — crescimento exponencial; investimentos, financiamentos
\end{itemize}
}

\nota{\textbf{Atenção:} no Brasil, a maioria das operações financeiras (CDB, poupança, financiamentos, cartões) utiliza o regime de \textbf{juros compostos}.}

\TW{Taxa}{badgepurple}{Taxas Proporcionais e Equivalentes}{Conversão de taxas entre períodos diferentes}

\regra{
\textbf{Taxas proporcionais} — regime simples; conversão por regra de três:
\[
  i_2 = i_1 \cdot \frac{t_2}{t_1}
\]

\textbf{Taxas equivalentes} — regime composto; conversão por potenciação:
\[
  (1 + i_a) = (1 + i_m)^{12}
\]
}

\exemp{
\textbf{Taxa mensal 3\% → anual proporcional (JS):}

$i_a = 3\% \cdot 12 = \hlm{36\%\text{ a.a.}}$

\textbf{Taxa anual 26,82\% → mensal equivalente (JC):}

$(1 + 0{,}2682) = (1 + i_m)^{12}$

$i_m = (1{,}2682)^{1/12} - 1 \approx 0{,}02$

$\hlm{i_m \approx 2\%\text{ ao mês}}$
}

\nota{\textbf{Atenção:} nunca converta taxas de juros compostos usando regra de três — isso gera erro. Use sempre a fórmula de equivalência.}

\TW{DS}{darkblue}{Desconto Simples}{Redução linear sobre o valor nominal}

\regra{
O \hl{desconto simples} calcula a redução \textbf{linear} no valor nominal de um título. O desconto incide sempre sobre o valor de face.
}

\formula{$D = N \cdot i \cdot t$}

\begin{itemize}
  \item $D$ — desconto
  \item $N$ — valor nominal do título
  \item $i$ — taxa de desconto por período (em decimal)
  \item $t$ — prazo até o vencimento
\end{itemize}

\SH{Valor Atual}

\formula{$A = N - D = N \cdot (1 - i \cdot t)$}

\exemp{
\textbf{Título R\$ 1.000,00, desconto 10\% a.a., prazo 1 ano:}

$D = 1000 \cdot 0{,}10 \cdot 1 = 100$

$A = 1000 - 100$

$\hlm{A = \text{R\$ } 900{,}00}$
}

\TW{DC}{darkblue}{Desconto Composto}{Desconto aplicado sucessivamente a cada período}

\regra{
No \hl{desconto composto}, o desconto é aplicado de forma \textbf{sucessiva} sobre o saldo remanescente de cada período — é o inverso do regime de juros compostos.
}

\formula{$A = N \cdot (1 - i)^t$}

O desconto total: $D = N - A$

\exemp{
\textbf{Título R\$ 5.000,00, desconto composto 5\% a.m. por 3 meses:}

$A = 5000 \cdot (0{,}95)^3 = 5000 \cdot 0{,}857375$

$A \approx \text{R\$ } 4.287{,}00$

$\hlm{D = 5000 - 4287 = \text{R\$ } 713{,}00}$
}

\TW{D∥D}{badgepurple}{Desconto por Fora vs. por Dentro}{Desconto comercial e desconto racional}

\regra{
\textbf{Desconto comercial (por fora):} calculado sobre o valor nominal $N$.
\[
  D_{\text{fora}} = N \cdot i \cdot t \quad A = N(1 - it)
\]

\textbf{Desconto racional (por dentro):} calculado sobre o valor atual $A$.
\[
  D_{\text{dentro}} = A \cdot i \cdot t \quad A = \frac{N}{1 + it}
\]
}

\exemp{
\textbf{Título R\$ 1.000,00, 10\% a.a., 1 ano:}

Por fora: $A = 1000 \cdot (1 - 0{,}10) = \text{R\$ } 900{,}00$

Por dentro: $A = \dfrac{1000}{1 + 0{,}10} \approx \text{R\$ } 909{,}09$

\hl{Desconto racional é sempre menor que o comercial.}
}

\nota{\textbf{Regra:} desconto comercial (por fora) favorece o banco; desconto racional (por dentro) favorece o devedor.}

\TW{Prob}{badgepurple}{Problemas Contextualizados}{Aplicando juros e desconto em situações reais}

\exemp{
\textbf{Prob. 1 — Juros Simples:}

R\$ 8.000,00 a 4\% a.m. por 5 meses.

$J = 8000 \cdot 0{,}04 \cdot 5 = 1600$

$\hlm{M = \text{R\$ } 9.600{,}00}$

\textbf{Prob. 2 — Juros Compostos:}

R\$ 2.000,00 a 12\% a.a. por 3 anos.

$M = 2000 \cdot (1{,}12)^3 = 2000 \cdot 1{,}404928$

$M \approx \text{R\$ } 2.809{,}86$

$\hlm{J \approx \text{R\$ } 809{,}86}$

\textbf{Prob. 3 — Desconto Comercial:}

Duplicata R\$ 3.500,00, 90 dias, taxa 3\% a.m.

$90\,\text{dias} = 3\,\text{meses}$

$D = 3500 \cdot 0{,}03 \cdot 3 = 315$

$\hlm{A = \text{R\$ } 3.185{,}00}$
}

\end{multicols}

\end{document}
